% !Mode:: "TeX:UTF-8"
% !Mode:: "TeX:UTF-8"
\newcommand{\thepage}{} % to silence a hyperref warning about pdfpagelabels
\RequirePackage[]{silence}
\WarningsOff[hyperref]
\PassOptionsToPackage{log-declarations=false}{xparse}
%\RequirePackage[l2tabu, orthodox]{nag}
\RequirePackage{luatex85} % to make \pdf...= commands work with LuaTeX
\documentclass[aspectratio=169,t]{beamer}
\usepackage{ifthen}
\usepackage[T1]{fontenc}
\usepackage{iftex}
\ifLuaTeX
 \usepackage{fontspec}
 \setmainfont[Ligatures=TeX]{Source Sans Pro}
\else
 \usepackage[utf8]{inputenc}
\fi
\usepackage{german}
\usepackage{newunicodechar}
\newunicodechar{€}{\texteuro}
\newunicodechar{ß}{\ss}
\renewcommand{\SS}{\iffontchar\font"1E9E \symbol{"1E9E}\else SS\fi}
\usepackage{ellipsis}
\usepackage{microtype}
\usepackage{listings}
\usepackage{sourcecodepro}
\lstset{language=Python, basicstyle=\ttfamily}
\usepackage{amsfonts}
\usepackage{amssymb}
\usepackage{amsthm}
\usepackage{fancyhdr}
\usepackage{graphicx}

\newtheorem{proposition}[theorem]{Proposition}

\DeclareMathOperator*{\argmax}{argmax}
\DeclareMathOperator{\supp}{supp}
\newcommand{\E}{\ensuremath{\mathbb{E}}}
\newcommand{\V}{\ensuremath{\mathbb{V}}}

\setbeamercolor{emph}{fg=purple}
\renewcommand<>{\emph}[1]{%
  {\usebeamercolor[fg]{emph}\only#2{\itshape}#1}%
}



\hypersetup{draft=false,colorlinks=false,pdfpagemode=UseOutlines,plainpages=false,
            breaklinks=true,bookmarksnumbered=true,pdfborder=0 0 0,
            pdftitle={Data Science SoSe 2026 HHU},
            pdfauthor={Konrad Voelkel},
            pdfsubject={data science}}

%\documentclass[aspectratio=169,t]{beamer}
\usetheme{hhubeamer} % put the hhu-beamer-theme in /usr/share/texmf/tex/latex/beamer/themes/

% -------------------- Code listings config --------------------
% Define custom colors for syntax highlighting
\definecolor{codegreen}{rgb}{0,0.6,0}
\definecolor{codegray}{rgb}{0.5,0.5,0.5}
\definecolor{codepurple}{rgb}{0.58,0,0.82}
\definecolor{backcolour}{rgb}{0.96,0.96,0.94} % Very light gray/beige
% Configure the listings package
\lstset{
    language=Python,
    backgroundcolor=\color{backcolour},   
    commentstyle=\color{codegreen}\itshape,
    keywordstyle=\color{blue}\bfseries,
    numberstyle=\tiny\color{codegray},
    stringstyle=\color{codepurple},
    basicstyle=\ttfamily\footnotesize, % Use a monospaced font
    breakatwhitespace=false,         
    breaklines=true,      
             % Enable automatic line breaking
    captionpos=b,                      % Caption at the bottom
    keepspaces=true,                 
    numbers=left,                      % Line numbers on the left
    numbersep=10pt,                    % Distance of line numbers from the code
    showspaces=false,                
    showstringspaces=false,            % Don't show weird underscores in strings
    showtabs=false,                  
    tabsize=4,
    frame=single,                      % Add a subtle border around the code
    rulecolor=\color{black!20},        % Border color
  }

\event[]{}
\author[Sommersemester 2026]{Dr. Konrad Völkel, Sommersemester 2026}

%\title{Data Science, VL 1}
%\date{16.04.2026}


%%% Local Variables:
%%% mode: latex
%%% TeX-master: t
%%% End:


\title{Data Science, VL 16}
\subtitle{Maximum Likelihood}
\date{18.06.2026}

\begin{document}

\begin{frame}
  \titlepage
\end{frame}

\section{Maximum Likelihood}

\begin{frame}[separator]
  \frametitle{Das Maximum-Likelihood-\\Prinzip}
\end{frame}
  
  
      \begin{frame}
    \frametitle{Likelihood-Funktion}
\begin{definition}
Gegeben ein statistisches Standardmodell $(\mathcal{X}, P_\theta : \theta \in \Theta)$ mit Dichte $f_\theta$ (entweder pmf oder pdf) für $P_\theta$ heißt $\rho \colon \mathcal{X} \times \Theta \to \mathbb{R}_{\geq 0}$ mit $\rho(x, \theta) := L_x(\theta) := f_\theta(x)$ die \alert{Likelihood-Funktion} und für festes $x$ heißt $L_x =\rho(x, \cdot) \colon \Theta \to \mathbb{R}_{\geq 0}$ die \alert{Likelihood}-Funktion zum Beobachtungswert $x$.
\end{definition}

\uncover<2->{
Erinnerung: Wenn wir die Hypothese $H = ( \theta$ ist der wahre Parameter$)$ haben, und $D=x$ (für Daten) schreiben, so ist $P(D|H) = L_x(\theta) = f_\theta(x)$ die Likelihood.}
  \end{frame}

  
      \begin{frame}
    \frametitle{MLE}
\begin{definition}
  Ein Schätzer $T \colon \mathcal{X} \to \Theta$ für den Parameter $\theta$ heißt \alert{Maximum-Likelihood-Schätzer} (maximum likelihood estimator, MLE), wenn für alle $x \in \mathcal{X}$ gilt:
\[L_x(T(x)) = \max_{\theta \in \Theta} L_x(\theta),\]
also wenn $T(x)$ die Likelihood von $x$ maximiert, für alle $x$. Man schreibt auch
\[T(x) = \argmax_{\theta \in \Theta} L_x(\theta).\]
\end{definition}
\end{frame}

  
      \begin{frame}
    \frametitle{Beispiel}
Im $N$-fachen univariat normalverteilten Modell ist $x \mapsto (\texttt{mean}(x), \sigma(x))$ ein Maximum-Likelihood-Schätzer für den Parameter $(\mu(X), \sigma(X))$. Der Schätzer $\sigma^\ast$ ist hingegen kein Maximum-Likelihood-Schätzer für $\sigma(X)$. Also können sich Erwartungstreue und Maximum-Likelihood-Prinzip gegenseitig ausschließen.
  \end{frame}

  
      \begin{frame}
    \frametitle{Bemerkung} 
    Mit Hilfe der relativen Entropie lässt sich beweisen, dass unter schwachen Voraussetzungen eine Folge von Maximum-Likelihood-Schätzern stets konsistent ist.

    \uncover<2->{
      Erinnerung: das heisst dass die Schätzwerte gegen den zu schätzenden Wert konvergieren.
      }
  \end{frame}

  
      \begin{frame}
        \frametitle{Bemerkung}
        \begin{itemize}[<+->]
          \item
    Wir sind also für eine Familie von Verteilungen $P_\theta$ mit Dichten $f_\theta$ daran interessiert, die Likelihood $f_\theta(x)$ einer Beobachtung $x$ in $\theta$ zu maximieren.
          \item
    Um Funktionen zu maximieren, betrachten wir gern die Ableitung\\ (dazu sollte $f$ in $\theta$ differenzierbar sein).
          \item
            Schon am Beispiel der univariaten Normalverteilungsdichte $\phi_{\mu,\sigma}$ sehen wir, dass die Ableitung $\frac{d}{d\mu} \phi_{\mu,\sigma}(x) =\frac{(x-\mu)}{\sigma^2} \phi_{\mu,\sigma}(x)$ bei $x=\mu$ verschwindet,\\
            damit ist $\mu=x$ der MLE für eine Stichprobe der Größe $1$.
          \item
            Im Allgemeinen ist die Likelihood allerdings schwerer zu berechnen \\
            (betrachte ein $n$-faches Produktmodell).
        \end{itemize}
      \end{frame}

  
      \begin{frame}
    \frametitle{Trick}
    Anstatt die Likelihood von $x$ über alle möglichen $\theta$ zu maximieren, kann man auch die \alert{log-Likelihood} $\log \rho(x,\theta)$ maximieren.

    $\log$ konkav $\Rightarrow$ ein Maximum der log-Likelihood ist auch ein Max. der Likelihood.

\uncover<2->{
Wenn z.B. $\rho$ die Form $e^{A}e^{B}$ hat (etwa bei einem $2$-fachen Produktmodell), hat $\log\rho$ die Form $A+B$ und Ableitungen gleich $0$ setzen wird deutlich einfacher.}
  \end{frame}

  
      \begin{frame}
        \frametitle{Beispiel Server-Logfiles}
        \begin{itemize}[<+->]
          \item
    Wir betrachten Server-Logfiles einer Web-Anwendung, insbesondere die Zeitstempel. Daraus können wir die konsekutiven Abstände berechnen, also die Wartezeit zwischen zwei Aufrufen.

    Wir möchten gerne wissen, mit welcher Wartezeit wir im Mittel rechnen können (um zusammen mit der durchschnittlichen Laufzeit eines Aufrufs zu einer Einschätzung zu kommen, welche Last wir erwarten können).

\item
      Die Wartezeiten sind exponentialverteilt mit Parameter $\lambda > 0$, den wir schätzen wollen. Zur Erinnerung: Die Dichte einer Exponentialverteilung ist $f(x) = \lambda e^{-\lambda x}$ und der Erwartungswert ist $\frac{1}{\lambda}$.
\end{itemize}
  \end{frame}
      
      \begin{frame}
        \frametitle{Beispiel Server-Logfiles}
    Wir betrachten Server-Logfiles (Zeitstempel) und konsekutive Abstände (Wartezeiten zwischen zwei Aufrufen). Die Wartezeiten sind exponentialverteilt mit Parameter $\lambda > 0$, den wir schätzen wollen. Die Dichte ist $f(x) = \lambda e^{-\lambda x}$.

Die Likelihood für eine Stichprobe der Größe $N$ ist $L_x(\lambda) = \prod_{i=1}^N \lambda e^{-\lambda x_i}$ und die log-Likelihood ist demnach

$$
\ell_x(\lambda) := \log L_x(\lambda) = \left( \sum_{i=1}^N \log\lambda - \lambda x_i\right) = n \left( \log\lambda - \lambda \texttt{mean}(x) \right).
$$


  \end{frame}
      
      \begin{frame}
    \frametitle{Beispiel Server-Logfiles (pt. 2)}
$$
\ell_x(\lambda) := \log L_x(\lambda) = \left( \sum_{i=1}^N \log\lambda - \lambda x_i\right) = n \left( \log\lambda - \lambda \texttt{mean}(x) \right).
$$

Die Ableitung nach $\lambda$ ist einfach $\frac{n}{\lambda} - n \texttt{mean}(x)$, eine Nullstelle liegt bei $\hat{\lambda} = \frac{1}{\texttt{mean}(x)}$. Die zweite Ableitung von $\ell_x$ ist $-\frac{n}{\lambda^2}$, was bei $\hat{\lambda}$ den Wert $-n{\texttt{mean}(x)}^2 < 0$ annimmt, also ist $\hat{\lambda}$ ein Maximum und wir haben den Maximum-Likelihood-Schätzer bestimmt.

\vspace{1cm}

\uncover<2->{
  Wir rechnen also mit einer erwarteten Wartezeit $\texttt{mean}(x)$ zwischen den Server-Aufrufen, wobei $x$ die beobachteten Wartezeiten sind.
}
  \end{frame}

      \begin{frame}
    \frametitle{Beispiel k-Means} 
Der k-Means-Algorithmus kann als Schätzer interpretiert werden:
das statistische Modell ist $\mathcal{X} = \mathbb{R}^{nN}$, wenn $n$ die Dimension der Daten ist (also bei $4$ reellen Variablen, etwa im Iris-Datensatz, $n=4$) und $N$ die Stichprobengröße (bei Iris ist $N=150$).

Der Parameterraum ist $\mathbb{R}^{nk}$ und die Verteilung zu einem Parameter $\theta = (c_1,\dots,c_k) \in \mathbb{R}^{nk}$ (mit $c_i \in \mathbb{R}^n$) ist ein gemischtes Modell aus multivariaten Normalverteilungen mit Erwartungswert $c_i$ und der $n\times n$-Einheitsmatrix als Kovarianzmatrix.
  \end{frame}

  
      \begin{frame}
    \frametitle{Beispiel k-Means} 
Die 'perfekte' k-Means-Lösung, also die Zentroide $(c_1,\dots,c_k)$, die der k-Means-Algorithmus versucht, zu bestimmen, sind diejenigen, für die die Summe der euklidischen Abstände eines jeden Punktes $x_i$ eines Samples $x \in \mathcal{X}$ minimal sind. Das sind genau die Zentroide, für die die Likelihood der Daten unter der entsprechenden gemischten Verteilung maximal wird.
\end{frame}

  
%\begin{frame}[separator]
%  \frametitle{Vielen Dank für Ihre Aufmerksamkeit!}
%\end{frame}

\end{document}

%%% Local Variables:
%%% mode: latex
%%% TeX-master: t
%%% End:
